\cleardoublepage

\chapter{Wstęp i cel pracy}
Od 2022 roku duże modele językowe stały się nieodłącznym narzędziem w życiu każdego człowieka, ułatwiając codzienne zadania. Każdego roku ukazują się nowe, bardziej zaawansowane i przełomowe LLMy. Obecnym trendem w świecie biznesu jest zastosowanie komercyjnych modeli, zapewniając im dostęp do wewnętrznej dokumentacji, baz danych czy repozytoriów kodu. W takiej sytuacji użytkownik lub firma stoi przed wyborem optymalnego rozwiązania, które jest możliwie skuteczne, tanie i szybko działające. Stąd, w tym dynamicznym środowisku, wychodzi potrzeba cyklicznego testowania i porównywania najnowszych popularnych algorytmów biorąc pod uwagę powyższe kryteria.


    \section{Cel pracy}
    Celem pracy magisterskiej jest stworzenie projektu, analiza wydajności oraz implementacja rozwiązania, umożliwiającego ekstrakcję i agregację danych z ustrukturyzowanych i nieustrukturyzowanych źródeł informacji. W ramach powyższych aktywności należy wykorzystać podejście Retrieval-Augmented Generation (RAG) współpracujące z dużymi modelami językowymi. W ramach ewaluacji skuteczności rozwiązania wskazane jest  przeprowadzenie testów porównawczych różnych wariantów systemów RAG, które obejmują m.in.: techniki podziału dokumentów (chunking), różne embeddery, zastosowanie wektorowych baz danych oraz wybranych popularnych najnowszych modeli językowych. W zakresie analizy skuteczności wskazane jest przeprowadzenie oceny w obszarze trafności odpowiedzi, kosztu generacji (tokenów), czasu odpowiedzi, odporności na pytania spoza kontekstu oraz zdolności algorytmów do odmowy udzielenia odpowiedzi w przypadku braku wystarczającej ilości danych. Najlepiej dobrana konfiguracja RAG powinna zostać udostępniona poprzez graficzny interfejs użytkownika, umożliwiając użytkownikowi zadawanie pytań opartych na dostępie do własnych źródeł danych. Opracowane rozwiązanie powinno wspierać generowanie spersonalizowanych rekomendacji dla odbiorców zainteresowanych konkretnymi aspektami usług lub informacji w wybranym obszarze tematycznym.